\section{Introdução}

Este trabalho tem como objetivo analisar aspectos da dinâmica demográfica do Amapá, considerando o período de 2000 a 2024. O recorte territorial adotado foi a Unidade da Federação Amapá, identificada pelo código 16 do IBGE. A escolha desse recorte permitiu trabalhar, de forma integrada, informações de nascimentos, óbitos, população residente e estrutura por idade e sexo.

Foram utilizados principalmente os dados do SINASC e do SIM, disponibilizados pelo DATASUS, além das projeções populacionais do IBGE -- Revisão 2024 e dos dados do Censo Demográfico de 2022. A partir dessas bases, foram calculados indicadores de sobrevivência na infância, natalidade, fecundidade, reprodução, mortalidade e tábuas de vida.

O trabalho está organizado em três partes principais. A primeira trata do Diagrama de Lexis e da sobrevivência na infância. A segunda analisa os indicadores de natalidade e fecundidade. A terceira apresenta os indicadores de mortalidade, a estrutura de causas de morte e as tábuas de vida. Ao longo das análises, buscou-se manter atenção à qualidade dos dados, aos denominadores utilizados e às limitações dos registros administrativos.
